\chapter{Literature Review}
\label{ch:literature_review}

\section{Introduction}
Understanding stock return predictability requires tracing the evolution of empirical asset pricing over the past six decades. Early academic frameworks established that asset prices in competitive markets reflect available information \citep{fama1970efficient}, modeling expected excess returns as linear functions of systematic risk exposures \citep{sharpe1964capital, lintner1965valuation}. Subsequent econometric advances operationalized these principles through multifactor cross-sectional regression models \citep{fama1973risk, ross1976arbitrage, fama1993common, fama2015five}.

However, the rapid expansion of financial data led to what \citet{cochrane2011discount} famously termed the ``factor zoo'', a proliferation of over 300 candidate anomaly variables claiming to generate alpha. In response, modern financial econometrics has developed along two primary paths: (i) statistical auditing of data-mining inflation and false discovery rates \citep{harvey2016and, harvey2020false, hou2020replicating}, and (ii) high dimensional regularized estimation and nonparametric machine learning \citep{kozak2020shrinking, freyberger2020dissecting, gu2020empirical, chen2023deep}.

In this chapter, I synthesize the asset pricing literature around six key thematic pillars that ground the theoretical framework, model architectures, feature selection, and statistical inference used in this thesis:
\begin{enumerate}
    \item \textbf{Theoretical Asset Pricing Foundations}: From the single factor equilibrium of CAPM and APT to the multifactor models of Fama-French, Carhart, and Pastor-Stambaugh.
    \item \textbf{The Factor Zoo, Multiple Testing, and Shrinkage}: Multiple testing inflation, false discovery rates, and high dimensional regularized SDF estimation.
    \item \textbf{High-Dimensional SDF Estimation \& Characteristic Interactions}: Nonparametric characteristic selection, latent factor models, and dynamic risk loadings.
    \item \textbf{Machine Learning \& Deep Learning Paradigms}: Nonlinear empirical asset pricing, instrumented principal components, and deep sequence SDF models.
    \item \textbf{Financial Time-Series Econometrics \& Inference}: HAC covariance estimation, out-of-sample forecast accuracy tests, and data-mining diagnostics.
    \item \textbf{Financial ML Protocols \& Overfitting Mitigation}: Backtesting rules, expanding walk-forward splits, microstructural frictions, and transaction cost frameworks.
\end{enumerate}

\section{Theoretical Asset Pricing Foundations}

\subsection{Equilibrium Foundations: The CAPM and Arbitrage Pricing Theory}
Modern quantitative finance begins with Markowitz's \citeyearpar{markowitz1952portfolio} mean-variance framework and the Capital Asset Pricing Model (CAPM) of \citet{sharpe1964capital} and \citet{lintner1965valuation}. CAPM assumes homogeneous investor expectations, frictionless markets, and quadratic utility. Under these ideal conditions, every investor holds the market portfolio $R_m$, making it mean-variance efficient. An asset's expected excess return is governed strictly by its market beta:
\begin{equation}
\mathbb{E}[R_{i,t}] - R_{f,t} = \beta_{i,m} \left( \mathbb{E}[R_{m,t}] - R_{f,t} \right)
\label{eq:ch2_capm}
\end{equation}
where $\beta_{i,m} = \frac{\operatorname{Cov}(R_{i,t}, R_{m,t})}{\operatorname{Var}(R_{m,t})}$ captures systemic market risk. \citet{black1972capital} relaxed risk-free borrowing assumptions, introducing the zero-beta CAPM. Black showed that borrowing constraints flatten the Security Market Line while raising its intercept.

To address CAPM's equilibrium restrictions, \citet{ross1976arbitrage} proposed Arbitrage Pricing Theory (APT). Grounded in asymptotic arbitrage-free pricing rather than market equilibrium, APT specifies a linear $K$-factor generative return structure:
\begin{align}
\label{eq:ch2_apt}
R_{i,t} &= \mathbb{E}[R_{i,t}] + \sum_{k=1}^{K} \beta_{i,k}\, f_{k,t} + \epsilon_{i,t}
\end{align}
where $f_{k,t}$ represent mean-zero common macroeconomic risk factors, $\beta_{i,k}$ are factor loadings, and $\epsilon_{i,t}$ is idiosyncratic noise with $\operatorname{Cov}(\epsilon_{i,t},\,\epsilon_{j,t})=0$ for $i \neq j$. Applying linear pricing space arguments in Hilbert spaces, Ross established that expected returns are linearly spanned by factor risk premia $\lambda_k$:
\begin{equation}
\mathbb{E}[R_{i,t}] - R_{f,t} = \sum_{k=1}^{K} \beta_{i,k} \lambda_k
\label{eq:ch2_apt_premia}
\end{equation}

\subsection{The Fama-MacBeth Two-Pass Regression Framework}
To test whether factor loadings explain cross-sectional stock returns, \citet{fama1973risk} introduced a two-pass regression methodology. In Pass 1, I run time-series regressions asset by asset over a rolling historical lookback window $T_1$ to estimate factor betas $\hat{\beta}_{i,k}$:
\begin{equation}
R_{i,t} - R_{f,t} = \alpha_i + \sum_{k=1}^{K} \beta_{i,k} f_{k,t} + \epsilon_{i,t}, \quad t = 1, \dots, T_1
\label{eq:ch2_fm_first_pass}
\end{equation}
In Pass 2, I run daily cross-sectional regressions across all available equities $i=1,\dots,N$ using estimated factor betas as explanatory variables:
\begin{equation}
R_{i,t} - R_{f,t} = \gamma_{0,t} + \sum_{k=1}^{K} \hat{\beta}_{i,k} \gamma_{k,t} + u_{i,t}, \quad i = 1, \dots, N
\label{eq:ch2_fm_second_pass}
\end{equation}
The point estimates of factor risk premia $\bar{\gamma}_k$ and their standard errors are obtained by averaging cross-sectional slope coefficients over time:
\begin{equation}
\bar{\gamma}_k = \frac{1}{T_2} \sum_{t=1}^{T_2} \hat{\gamma}_{k,t}, \quad \text{SE}(\bar{\gamma}_k) = \sqrt{\frac{1}{T_2(T_2-1)} \sum_{t=1}^{T_2} \left( \hat{\gamma}_{k,t} - \bar{\gamma}_k \right)^2}
\label{eq:ch2_fm_se}
\end{equation}
Because standard Fama-MacBeth standard errors ignore serial correlation and generated-regressor error, modern applications apply Heteroskedasticity and Autocorrelation Consistent (HAC) corrections \citep{newey1987simple} and \citet{shanken1992on} beta-uncertainty adjustments.

\subsection{Multi-Factor Models: Fama-French, Carhart, and Pastor-Stambaugh}
Single-factor CAPM faced empirical challenges in the 1980s as market anomalies accumulated. \citet{fama1992cross} documented that size (market capitalization) and book-to-market ($B/M$) ratios absorbed much of market beta's explanatory power. To capture these anomalies, \citet{fama1993common} formulated the Fama-French Three-Factor Model:
\begin{equation}
R_{i,t} - R_{f,t} = \alpha_i + \beta_{i,1} (R_{m,t} - R_{f,t}) + \beta_{i,2} SMB_t + \beta_{i,3} HML_t + e_{i,t}
\label{eq:ch2_ff3}
\end{equation}
where $SMB_t$ (Small Minus Big) isolates the size premium, and $HML_t$ (High Minus Low) isolates the value premium.

To address the price momentum anomaly \citep{jegadeesh1993returns}, \citet{carhart1997on} added a momentum factor ($WML_t$, Winners Minus Losers):
\begin{equation}
R_{i,t} - R_{f,t} = \alpha_i + \beta_{i,1} (R_{m,t} - R_{f,t}) + \beta_{i,2} SMB_t + \beta_{i,3} HML_t + \beta_{i,4} WML_t + e_{i,t}
\label{eq:ch2_carhart4}
\end{equation}
Concurrently, \citet{pastor2003liquidity} demonstrated that aggregate market liquidity represents an un-diversifiable risk factor. Stocks with high liquidity sensitivity command higher expected returns to compensate for illiquidity during market drawdowns.

To incorporate operating profitability and investment intensity, \citet{fama2015five} introduced their Five-Factor Asset Pricing Model:
\begin{align}
\label{eq:ch2_ff5_model}
R_{i,t} - R_{f,t} &= \alpha_i + \beta_{i,1} (R_{m,t} - R_{f,t}) + \beta_{i,2}\, SMB_t + \beta_{i,3}\, HML_t \nonumber \\
&\quad + \beta_{i,4}\, RMW_t + \beta_{i,5}\, CMA_t + e_{i,t}
\end{align}
where $RMW_t$ (Robust Minus Weak) measures profitability differentials, and $CMA_t$ (Conservative Minus Aggressive) measures investment intensity. Grounded in the dividend discount model, \citet{fama2015five} showed that holding book-to-market constant, higher profitability implies higher expected returns, while higher investment lowers expected returns. In parallel, \citet{hou2015digesting} introduced the $q$-factor model ($Mkt, ME, I/A, ROE$), establishing that investment and profitability factors derived from $q$-theory span a broad set of anomalies.

\section{The Factor Zoo Controversy, Multiple Testing, and Shrinkage}

\subsection{The Proliferation of Anomalies and Cochrane's "Factor Zoo"}
By 2011, academic papers claimed over 300 firm characteristics predicted stock returns. In his AFA presidential address, \citet{cochrane2011discount} dubbed this proliferation the ``factor zoo.'' Cochrane noted that while theoretical models assume a small set of macro factors span returns, empirical literature created dozens of redundant predictors without verifying if candidate factors were incremental to existing benchmarks.

This problem was worsened by post-publication return decay. \citet{mclean2016does} examined 97 published anomalies and showed that out-of-sample portfolio returns fall by 26\% after publication and up to 58\% post-sample. This decay reflects arbitrage capital exploiting published signals alongside data-mining bias in original studies.

\subsection{Multiple Testing Inflation and False Discovery Rates}
\citet{harvey2016and} analyzed the statistical cause of factor zoo inflation. Standard empirical studies tested candidate factors using single-hypothesis $t$-tests ($|t| > 1.96, \alpha = 0.05$). Evaluating hundreds of candidate signals on identical datasets creates high false discovery rates.

\citet{harvey2016and} evaluated 300+ anomaly factors under multiple-testing adjustments (Bonferroni, Holm-Bonferroni). They established that keeping the False Discovery Rate (FDR) below 5\% requires raising the minimum $t$-statistic threshold to:
\begin{equation}
|t| > 3.0 \quad (p \text{-value} < 0.0027)
\label{eq:ch2_harvey_threshold}
\end{equation}
\citet{harvey2020false} extended this via Bayesian FDR frameworks. In a large replication audit, \citet{hou2020replicating} re-evaluated 452 cross-sectional anomalies under microcap exclusions and value-weighting. They found that 65\% of published anomalies fail at $|t| > 1.96$, and over 85\% fail under $|t| > 3.0$.

\subsection{High-Dimensional Regularized SDF Estimation}
To address the factor zoo problem structurally, econometricians re-framed asset pricing as a high dimensional regularized Stochastic Discount Factor (SDF) estimation problem. In a stochastic discount factor framework, an asset pricing model specifies a scalar pricing kernel $M_{t+1}$ such that for all excess returns $R_{i,t+1}$:
\begin{equation}
\mathbb{E}_t \left[ M_{t+1} R_{i,t+1} \right] = 0
\label{eq:ch2_sdf_fundamental}
\end{equation}
When the SDF is parameterized as a linear combination of a high dimensional vector of $K$ candidate characteristics or factor returns $\mathbf{z}_{t+1}$, i.e., $M_{t+1} = 1 - \mathbf{b}^T (\mathbf{z}_{t+1} - \mathbb{E}[\mathbf{z}_{t+1}])$, unregularized sample estimators fail due to multicollinearity and $K \approx T$ sample noise.

\citet{kozak2020shrinking} showed that shrinking the factor weighting vector $\mathbf{b}$ via combined $L_1$ (LASSO) and $L_2$ (Ridge) regularization yields robust, out-of-sample optimal SDF estimators. The regularized SDF objective function is formulated as:
\begin{equation}
\min_{\mathbf{b}} \left\{ \frac{1}{T} \sum_{t=1}^{T} \left( 1 - \mathbf{b}^T \mathbf{z}_t \right)^2 + \gamma_1 \|\mathbf{b}\|_1 + \gamma_2 \|\mathbf{b}\|_2^2 \right\}
\label{eq:ch2_kozak_regularization}
\end{equation}
\citet{kozak2020shrinking} showed that while unregularized models tend to overfit high-frequency sample noise, penalized SDF estimation extracts a low-dimensional principal component structure that spans expected returns while shrinking noise coefficients to zero.

\section{High-Dimensional SDF Estimation \& Characteristic Interactions}

\subsection{Non-Parametric Characteristic Pruning via Group LASSO}
Parallel to regularized linear SDFs, \citet{freyberger2020dissecting} developed a nonparametric model to systematically select which firm characteristics provide incremental predictive power for expected returns. Using nonparametric B-spline expansions and Group LASSO penalties, their estimator models expected returns as:
\begin{equation}
\mathbb{E}_t[R_{i,t+1}] = \sum_{j=1}^{J} g_j(c_{i,j,t}) + \epsilon_{i,t+1}
\label{eq:ch2_freyberger_nonparametric}
\end{equation}
where $c_{i,j,t}$ represents the rank-transformed characteristic $j$ for firm $i$ at time $t$, and $g_j(\cdot)$ is a nonlinear continuous function. The Group LASSO penalty penalizes the vector of spline coefficients for each characteristic simultaneously:
\begin{equation}
\min_{\{g_j\}} \left\{ \frac{1}{NT} \sum_{i=1}^{N} \sum_{t=1}^{T} \left( R_{i,t+1} - \sum_{j=1}^{J} g_j(c_{i,j,t}) \right)^2 + \lambda \sum_{j=1}^{J} \sqrt{\text{dim}(g_j)} \|g_j\|_2 \right\}
\label{eq:ch2_group_lasso}
\end{equation}
\citet{freyberger2020dissecting} showed that out of more than 50 published characteristics, only a small subset (such as short term momentum, market capitalization, operating profitability, and asset growth) retain independent predictive power once non-linearities and characteristic redundancies are penalized.

\subsection{Sparse Signals and Latent Factor Shrinkage}
\citet{chinco2019sparse} examined sparse return predictability at high frequencies, showing that LASSO regression identifies transient, cross-asset predictive signals that standard time series estimators miss. Grounding model selection in econometric diagnostics, \citet{lewellen2010skeptical} provided a critique of cross-sectional asset pricing tests. They showed that testing multifactor models on strong factor-structure portfolios (such as standard 25 Size-B/M sorted portfolios) produces artificially inflated cross-sectional $R^2$ values. They prescribed rigorous diagnostic standards: expanding test portfolios beyond Size-B/M, forcing theoretical factor risk premia to equal sample factor means, and reporting confidence intervals for cross-sectional $R^2$.

Expanding structural factor selection, \citet{feng2020taming} and \citet{giglio2021test} developed double-testing econometric frameworks combining LASSO factor selection with two-pass regressions. They showed that high dimensional regularization enables researchers to evaluate whether a newly proposed factor adds incremental explanatory power to the existing universe of 300+ factors without suffering from omitted variable bias or data mining contamination.

\section{Machine Learning and Deep Learning Paradigms}

\subsection{Empirical Asset Pricing Benchmark: Gu, Kelly, and Xiu (2020)}
\citet{gu2020empirical} established the primary empirical benchmark for financial machine learning. Evaluating a 60-year panel of over 30,000 U.S. equities (1957--2016) with 94 firm characteristics and 8 macroeconomic indicators, Gu, Kelly, and Xiu tested linear OLS, LASSO, ElasticNet, PCR, PLS, Random Forests, Gradient Boosted Trees (GBDT), and Feedforward Neural Networks (MLP).

The non-linear forecasting panel model takes the form:
\begin{equation}
R_{i,t+1} = g^*(\mathbf{x}_{i,t}) + \epsilon_{i,t+1}, \quad \mathbf{x}_{i,t} = \mathbf{c}_{i,t} \otimes \mathbf{d}_t
\label{eq:ch2_gk_model}
\end{equation}
where $\mathbf{c}_{i,t}$ is a $p$-dimensional vector of firm characteristics, $\mathbf{d}_t$ is an $m$-dimensional vector of macro variables, and $\otimes$ denotes the Kronecker product creating interactive feature conditioning.

\citet{gu2020empirical} established three key empirical results:
\begin{enumerate}
    \item \textbf{Nonlinear Model Advantage}: Tree ensembles and neural networks outperform linear specifications in out-of-sample $R_{\text{oos}}^2$, effectively doubling long-short portfolio Sharpe ratios.
    \item \textbf{Characteristic Interaction Gains}: Most predictive gain stems from non-linear interaction terms (such as size interacted with momentum and liquidity).
    \item \textbf{Optimal Network Depth}: Shallow networks (1--3 hidden layers) outperform deeper architectures, as deeper networks overfit on low signal-to-noise equity return data.
\end{enumerate}

\subsection{Characteristics as Dynamic Risk Exposures: IPCA}
\citet{kelly2019characteristics} provided theoretical backing for why machine learning forecasts succeed. Introducing Instrumented Principal Component Analysis (IPCA), Kelly, Pruitt, and Su modeled asset returns as:
\begin{equation}
R_{i,t+1} = \mathbf{z}_{i,t}^T \mathbf{\Gamma}_{\beta} \mathbf{f}_{t+1} + \epsilon_{i,t+1}
\label{eq:ch2_ipca}
\end{equation}
where $\mathbf{z}_{i,t}$ is a vector of observable firm characteristics, $\mathbf{\Gamma}_{\beta}$ is a mapping matrix connecting characteristics to dynamic factor loadings ($\boldsymbol{\beta}_{i,t} = \mathbf{\Gamma}_{\beta}^T \mathbf{z}_{i,t}$), and $\mathbf{f}_{t+1}$ represents unobservable latent factor returns.

\citet{kelly2019characteristics} showed that firm characteristics proxy for time-varying factor risk exposures. Machine learning return predictors $R_{i,t+1} = g(\mathbf{z}_{i,t})$ are thus implicitly estimating dynamic latent factor models with time-varying risk loadings.

\subsection{Deep Learning and Non-Parametric SDF Estimation}
\citet{chen2023deep} extended machine learning to deep architectures, estimating non-parametric asset pricing SDFs directly from high-dimensional equity panels. Parameterizing factor loadings $\beta(\mathbf{c}_{i,t})$ via Deep Neural Networks (DNNs) combined with Generative Adversarial Networks (GANs), Chen, Pelger, and Zhu showed that deep learning captures macro regime shifts and structural non-linearities missed by linear factor models.

Similarly, \citet{cong2021textual} built multimodal deep learning models to process accounting signals alongside high-dimensional financial news text. To interpret deep neural network predictions, \citet{lundberg2017unified} introduced SHapley Additive exPlanations (SHAP), based on cooperative game theory:
\begin{equation}
f(\mathbf{x}) = \phi_0 + \sum_{j=1}^{M} \phi_j \mathbf{x}_j
\label{eq:ch2_shap}
\end{equation}
allowing transparent feature attribution in non-linear models.

\subsection{Economic Frictions and Limits to Arbitrage}
\citet{avramov2023machine} examined whether machine learning forecast gains survive real-world market frictions. Avramov, Cheng, and Metzker demonstrated that machine learning predictability is concentrated in illiquid, microcap stocks with high distress risk and short-sale constraints. Excluding microcaps or deducting realistic execution fees significantly reduces out-of-sample strategy profits. This highlights the necessity of incorporating transaction cost frictions into financial machine learning evaluation.

\section{Financial Time-Series Econometrics \& Inference}

\subsection{Heteroskedasticity and Autocorrelation Consistent (HAC) Inference}
Financial return series exhibit volatility clustering, heavy tails, and serial correlation. Standard OLS covariance matrices underestimate parameter variances, producing inflated $t$-statistics. To ensure valid statistical inference, \citet{newey1987simple} developed the Heteroskedasticity and Autocorrelation Consistent (HAC) covariance matrix estimator.

For regression residuals $e_t$ and regressors $\mathbf{X}_t$, defining $\mathbf{u}_t = \mathbf{X}_t e_t$, the Newey-West sample covariance matrix $\widehat{\Omega}_{\text{NW}}$ is:
\begin{equation}
\widehat{\Omega}_{\text{NW}} = \widehat{\Gamma}_0 + \sum_{j=1}^{L} w(j, L) \left( \widehat{\Gamma}_j + \widehat{\Gamma}_j^T \right)
\label{eq:ch2_newey_west_eq}
\end{equation}
where $\widehat{\Gamma}_j = \frac{1}{T} \sum_{t=j+1}^{T} \mathbf{u}_t \mathbf{u}_{t-j}^T$ is the sample autocovariance at lag $j$, and $w(j, L) = 1 - \frac{j}{L+1}$ is the Bartlett kernel weighting function over lag length $L$. The Bartlett kernel guarantees that $\widehat{\Omega}_{\text{NW}}$ remains positive semi-definite in finite samples.

\subsection{Out-of-Sample Predictive Accuracy: Diebold-Mariano and HLN Corrections}
When comparing out-of-sample forecast accuracy between competing models (e.g. TFDMGA versus PyTorch LSTM), standard independent sample $t$-tests fail due to forecast error serial dependence. \citet{diebold1995comparing} introduced a distribution-free statistic for testing equal predictive accuracy.

Let $e_{1,t}$ and $e_{2,t}$ denote forecast errors from Model 1 and Model 2 at time $t$. The loss differential sequence under Mean Squared Error is:
\begin{equation}
d_t = L(e_{1,t}) - L(e_{2,t})
\label{eq:ch2_loss_diff}
\end{equation}
Under the null hypothesis $H_0: \mathbb{E}[d_t] = 0$, the Diebold-Mariano statistic is:
\begin{equation}
DM = \frac{\bar{d}}{\sqrt{\frac{\widehat{V}(\bar{d})}{T}}} \xrightarrow{d} \mathcal{N}(0, 1)
\label{eq:ch2_dm_stat}
\end{equation}
where $\bar{d} = \frac{1}{T}\sum_{t=1}^T d_t$, and $\widehat{V}(\bar{d})$ is the Newey-West HAC-adjusted asymptotic variance.

To correct for small-sample size distortion, \citet{harvey1997testing} introduced the Harvey-Leybourne-Newbold (HLN) correction. For forecast horizon $h$, the HLN-adjusted statistic is:
\begin{equation}
DM^* = DM \times \left[ \frac{T + 1 - 2h + T^{-1}h(h-1)}{T} \right]^{1/2}
\label{eq:ch2_hln_stat}
\end{equation}
Critical values for $DM^*$ are evaluated against the Student's $t$-distribution with $(T-1)$ degrees of freedom.

\subsection{Data Mining Diagnostics: Reality Check and Superior Predictive Ability}
When tuning multiple model architectures across hyperparameters, researchers risk selecting a model that performs well purely by chance. To control for data-mining bias, \citet{white2000reality} developed the Reality Check for Data Mining using a bootstrap procedure over candidate model forecasts.

\citet{hansen2005superior} refined this framework into the Superior Predictive Ability (SPA) test, introducing a studentized test statistic with sample-dependent variance adjustments:
\begin{equation}
T_k^{\text{SPA}} = \max_{k=1, \dots, K} \frac{\sqrt{T} \bar{d}_k}{\hat{\omega}_k}
\label{eq:ch2_hansen_spa}
\end{equation}
where $\hat{\omega}_k^2$ estimates the asymptotic variance of $\sqrt{T}\bar{d}_k$. Hansen showed that SPA provides higher power and robustness against uninformative candidate models.

\section{Financial ML Protocols \& Overfitting Mitigation}

\subsection{Methodological Pitfalls in Financial Machine Learning}
In financial machine learning, standard statistical cross-validation techniques (such as standard randomized $K$-fold CV) can fail. As detailed by \citet{lopez2018advances}, financial time series data violate the independent and identically distributed (i.i.d.) assumption. Applying naive cross-validation induces several methodological problems:
\begin{enumerate}
    \item \textbf{Lookahead Leakage}: Using future information (such as future target returns or lookahead normalized features) during training.
    \item \textbf{Serial Memory Contamination}: Overlapping target return windows (such as multi-day cumulative returns) allow training samples to overlap with test samples, leaking labels across folds.
    \item \textbf{Backtest Overfitting}: Iteratively tuning model parameters directly on test set Sharpe ratios until an illusion of profitability is produced.
\end{enumerate}

\subsection{Validation Protocols}
To ensure backtesting integrity, \citet{lopez2018advances} established validation protocols:
\begin{itemize}
    \item \textbf{Expanding Walk-Forward Validation}: Models are trained on historical data $t \in [1, T_{\text{train}}]$ and evaluated on subsequent out-of-sample periods $t \in [T_{\text{train}}+1, T_{\text{test}}]$. The training window expands iteratively over time.
    \item \textbf{Combinatorial Purged Cross-Validation (CPCV)}: When non-sequential cross-validation is conducted, training folds adjacent to test folds are \emph{purged} by removing samples whose label windows overlap with test windows. In addition, an \emph{embargo} period is appended after test sets to eliminate serial correlation leakage.
    \item \textbf{Point-in-Time Data Alignment}: Financial fundamentals and accounting metrics are indexed by their actual public SEC disclosure dates (such as 10-K/10-Q filing dates) rather than fiscal period end dates.
\end{itemize}

\subsection{Microstructure Frictions \& Transaction Cost Taxonomy}
High-frequency predictive signals can decay rapidly once realistic transaction costs and market microstructure frictions are accounted for \citep{novy2016taxonomy, frazzini2014betting}. \citet{novy2016taxonomy} show that many published anomaly strategies disappear after deducting effective bid-ask spreads, market impact costs, and turnover frictions.

Second, execution realism mandates an execution delay: trading signals generated at market close on day $t$ cannot be executed until day $t+1$ close or open. In addition, leverage constraints and margin costs \citep{frazzini2014betting} restrict short-selling capacity on distressed equities, reinforcing the necessity of evaluating long-only and cost-adjusted long-short performance.

\section{Synthesis of Core Literature Benchmarks}
Table \ref{tab:ch2_lit_synthesis_comprehensive} provides a synthesis of the literature that anchors the theoretical framing, model architecture, feature engineering, econometric estimation, and validation of this thesis.

\begin{table}[htbp]
\centering
\caption{Comprehensive Synthesis of Core Literature Benchmarks}
\label{tab:ch2_lit_synthesis_comprehensive}
\resizebox{\textwidth}{!}{%
\begin{tabular}{llll}
\toprule
\textbf{Thematic Domain} & \textbf{Key Literature Citation} & \textbf{Core Theoretical / Econometric Finding} & \textbf{Implementation in Thesis} \\
\midrule
\multirow{4}{*}{\shortstack[l]{Asset Pricing\\Foundations}} 
& \citet{sharpe1964capital}, \citet{lintner1965valuation} & Single factor CAPM equilibrium; market beta linear pricing & Baseline risk benchmark \\
& \citet{fama1973risk} & Two-pass cross-sectional regression methodology & Baseline linear factor risk premia \\
& \citet{ross1976arbitrage} & Arbitrage Pricing Theory (APT); multifactor space & Linear multifactor baseline \\
& \citet{fama1993common}, \citet{fama2015five} & 3-Factor and 5-Factor models ($Mkt, SMB, HML, RMW, CMA$) & Econometric spanning regressions \\
& \citet{carhart1997on}, \citet{pastor2003liquidity} & Price momentum ($WML$) and aggregate market liquidity factors & Extended multifactor benchmark \\
\midrule
\multirow{4}{*}{\shortstack[l]{Factor Zoo \&\\Shrinkage}} 
& \citet{cochrane2011discount} & Factor zoo diagnosis; lack of unified factor structure & High dimensional model motivation \\
& \citet{harvey2016and}, \citet{harvey2020false} & Multiple testing bias; minimum hurdle rate $|t| > 3.0$ & Feature auditing threshold \\
& \citet{hou2020replicating}, \citet{mclean2016does} & 65\%+ anomaly replication failure; post-publication decay & Microcap exclusions \& value weighting \\
& \citet{kozak2020shrinking} & Regularized $L_1/L_2$ SDF shrinkage in high dimensions & Penalized Ridge/ElasticNet baselines \\
& \citet{freyberger2020dissecting} & Nonparametric Group LASSO characteristic pruning & Rank-normalized feature selection \\
\midrule
\multirow{4}{*}{\shortstack[l]{Machine Learning\\\& Deep Learning}} 
& \citet{gu2020empirical} & Nonlinear ML performance; interaction terms; 1--3 layer networks & Experimental lineup design \\
& \citet{kelly2019characteristics} & Instrumented PCA (IPCA); characteristics as dynamic betas & Latent dynamic factor framing \\
& \citet{chen2023deep} & Deep learning nonparametric SDF estimation via GANs & Sequence model architecture \\
& \citet{avramov2023machine} & ML gains concentrated in illiquid hard-to-arbitrage stocks & Size-stratified empirical evaluation \\
& \citet{lundberg2017unified} & SHAP additive feature attribution for nonlinear models & Model interpretability framework \\
\midrule
\multirow{4}{*}{\shortstack[l]{Time-Series\\Econometrics}} 
& \citet{newey1987simple} & HAC robust standard errors; Bartlett kernel weighting & Fama-MacBeth HAC t-statistics \\
& \citet{diebold1995comparing}, \citet{harvey1997testing} & Out-of-sample forecast accuracy test \& HLN correction & Model forecast comparison tests \\
& \citet{white2000reality}, \citet{hansen2005superior} & Reality Check and Superior Predictive Ability (SPA) tests & Hyperparameter data-mining audit \\
\midrule
\multirow{3}{*}{\shortstack[l]{Financial ML\\Validation}} 
& \citet{lopez2018advances} & Financial ML protocols; purging, embargoing, walk-forward & 5-Fold expanding walk-forward split \\
& \citet{novy2016taxonomy} & Transaction cost taxonomy; effective spreads and impact & 10 bps friction drag evaluation \\
& \citet{frazzini2014betting} & Betting against beta; leverage constraints and margin drag & Execution delay \& shorting rules \\
\bottomrule
\end{tabular}%
}
\end{table}

\section{Summary}
In conclusion, the evolution of empirical asset pricing shows a clear transition from simple linear factor models toward regularized high dimensional estimation and nonparametric deep sequence architectures. By synthesizing the theoretical insights of factor pricing, the statistical rigor of multiple testing adjustments, the predictive capability of deep sequence networks, and empirical validation protocols, this thesis builds a framework for cross-sectional stock return prediction.
